% GENERATED FILE. Edit canonical lessons, then run npm run generate:books.
% canonical-source-hash: fnv1a64:47dc2fca4f02821e
% canonical-lessons: JA-C280-kiji, JA-C280-noseru, JA-C280-dekigoto, JA-C280-kisha, JA-C280-kubaru

\chapter{What the News Said}
\label{ch:ja-what-the-news-said}
\hlchaptermodality{280}

\begin{chapteropening}
\textbf{By the end of this chapter:} \emph{I can say what the paper reported and who reported it.}

Complete the last lesson of chapter 280: I can say what the paper reported and who reported it.
\end{chapteropening}

\section[\texorpdfstring{\ja{きじ} (kiji)}{きじ (kiji)}]{\ja{きじ} (kiji) — a newspaper article; a report}

\label{lesson:JA-C280-kiji}



\begin{quote}
Before the new one: say the Japanese for to prevent, then the Japanese for big.
\end{quote}

\subsection*{You'll want to know: \ja{きじ}}
\textbf{\ja{きじ}} — \emph{kiji} — \textquotedblleft{}a newspaper article; a report\textquotedblright{}.

\begin{quote}
\textbf{\ja{今} \ja{しんぶんの} \ja{きじを} \ja{よんでいます。}} — \emph{ima shinbun no kiji o yonde imasu} — \textquotedblleft{}I am reading a newspaper article now.\textquotedblright{}
\end{quote}

\begin{cousinweb}
\textbf{\ja{きじ}} (\emph{kiji}) is Sino-Japanese (kango): two Chinese characters read with their on'yomi, readings borrowed from Chinese. They are \emph{ki}, \textquotedblleft{}account\textquotedblright{} (Middle Chinese \emph{kiH}) and \emph{ji}, \textquotedblleft{}matter\textquotedblright{} (\emph{dzriH}) (KANJIDIC2; Middle Chinese in Baxter's transcription).
\end{cousinweb}

\begin{grammarlens}[title={What you've built}]
One more word for telling what happened.
\end{grammarlens}

\subsection*{Your turn}
\begin{itemize}
\raggedright
  \item \emph{Say it:} \emph{kiji}
  \item \emph{Say it:} \emph{kiji}, once more
  \item \emph{Say it:} \emph{ima shinbun no kiji o yonde imasu} — I am reading a newspaper article now.
  \item \emph{Recall:} say the Japanese for to prevent, then the Japanese for big, then say \emph{kiji} again
\end{itemize}

\begin{tcolorbox}[breakable,colback=teal!4,colframe=teal!35!black,title={Before you move on}]
What does \ja{きじ} mean? (\textquotedblleft{}A newspaper article; a report\textquotedblright{}.) Say it once more.
\end{tcolorbox}
\section[\texorpdfstring{\ja{のせる} (noseru)}{のせる (noseru)}]{\ja{のせる} (noseru) — to print (in a paper); to put on}

\label{lesson:JA-C280-noseru}



\begin{quote}
Before the new one: say the Japanese for big, then the Japanese for a newspaper article.
\end{quote}

\subsection*{You'll want to know: \ja{のせる}}
\textbf{\ja{のせる}} — \emph{noseru} — \textquotedblleft{}to print (in a paper); to put on\textquotedblright{}.

\begin{quote}
\textbf{\ja{しんぶんが} \ja{じけんの} \ja{しゃしんを} \ja{のせました。}} — \emph{shinbun ga jiken no shashin o nosemashita} — \textquotedblleft{}The paper ran a photo of the incident.\textquotedblright{}
\end{quote}

Polite \textbf{\ja{のせます}}, \emph{nosemasu}; an ichidan verb. Past (polite): \textbf{\ja{しんぶんが} \ja{しゃしんを} \ja{のせました。}}, \emph{shinbun ga shashin o nosemashita}.

\begin{cousinweb}
\textbf{\ja{のせる}} (\emph{noseru}) is a native Japanese word (wago), inherited rather than borrowed from Chinese; its kanji was attached later as a kun'yomi (JMdict).
\end{cousinweb}

\begin{grammarlens}[title={What you've built}]
One more word for telling what happened.
\end{grammarlens}

\subsection*{Your turn}
\begin{itemize}
\raggedright
  \item \emph{Say it:} \emph{noseru}
  \item \emph{Say it:} \emph{noseru}, once more
  \item \emph{Say it:} \emph{shinbun ga jiken no shashin o nosemashita} — The paper ran a photo of the incident.
  \item \emph{Recall:} say the Japanese for big, then the Japanese for a newspaper article, then say \emph{noseru} again
\end{itemize}

\begin{tcolorbox}[breakable,colback=teal!4,colframe=teal!35!black,title={Before you move on}]
What does \ja{のせる} mean? (\textquotedblleft{}To print (in a paper); to put on\textquotedblright{}.) Say it once more.
\end{tcolorbox}
\section[\texorpdfstring{\ja{できごと} (dekigoto)}{できごと (dekigoto)}]{\ja{できごと} (dekigoto) — an event; a happening}

\label{lesson:JA-C280-dekigoto}



\begin{quote}
Before the new one: say the Japanese for a newspaper article, then the Japanese for to print (in a paper).
\end{quote}

\subsection*{You'll want to know: \ja{できごと}}
\textbf{\ja{できごと}} — \emph{dekigoto} — \textquotedblleft{}an event; a happening\textquotedblright{}.

\begin{quote}
\textbf{\ja{きのうの} \ja{できごとを} \ja{はなして} \ja{ください。}} — \emph{kinō no dekigoto o hanashite kudasai} — \textquotedblleft{}Please tell me what happened yesterday.\textquotedblright{}
\end{quote}

\begin{cousinweb}
\textbf{\ja{できごと}} (\emph{dekigoto}) is a native Japanese word (wago), inherited rather than borrowed from Chinese, built from \emph{deki}, \textquotedblleft{}coming about\textquotedblright{} (from \emph{dekiru}), and \emph{goto}, \textquotedblleft{}matter\textquotedblright{} (from \emph{koto}) (JMdict, KANJIDIC2).
\end{cousinweb}

\begin{grammarlens}[title={What you've built}]
One more word for telling what happened.
\end{grammarlens}

\subsection*{Your turn}
\begin{itemize}
\raggedright
  \item \emph{Say it:} \emph{dekigoto}
  \item \emph{Say it:} \emph{dekigoto}, once more
  \item \emph{Say it:} \emph{kinō no dekigoto o hanashite kudasai} — Please tell me what happened yesterday.
  \item \emph{Recall:} say the Japanese for a newspaper article, then the Japanese for to print (in a paper), then say \emph{dekigoto} again
\end{itemize}

\begin{tcolorbox}[breakable,colback=teal!4,colframe=teal!35!black,title={Before you move on}]
What does \ja{できごと} mean? (\textquotedblleft{}An event; a happening\textquotedblright{}.) Say it once more.
\end{tcolorbox}
\section[\texorpdfstring{\ja{きしゃ} (kisha)}{きしゃ (kisha)}]{\ja{きしゃ} (kisha) — a reporter; a journalist}

\label{lesson:JA-C280-kisha}



\begin{quote}
Before the new one: say the Japanese for to print (in a paper), then the Japanese for an event.
\end{quote}

\subsection*{You'll want to know: \ja{きしゃ}}
\textbf{\ja{きしゃ}} — \emph{kisha} — \textquotedblleft{}a reporter; a journalist\textquotedblright{}.

\begin{quote}
\textbf{\ja{きしゃが} \ja{けいさつに} \ja{しつもんしました。}} — \emph{kisha ga keisatsu ni shitsumon shimashita} — \textquotedblleft{}A reporter put questions to the police.\textquotedblright{}
\end{quote}

\begin{cousinweb}
\textbf{\ja{きしゃ}} (\emph{kisha}) is Sino-Japanese (kango): two Chinese characters read with their on'yomi, readings borrowed from Chinese. They are \emph{ki}, \textquotedblleft{}account\textquotedblright{} (Middle Chinese \emph{kiH}) and \emph{sha}, \textquotedblleft{}person\textquotedblright{} (\emph{tsyaeX}) (KANJIDIC2; Middle Chinese in Baxter's transcription).
\end{cousinweb}

\begin{grammarlens}[title={What you've built}]
One more word for telling what happened.
\end{grammarlens}

\subsection*{Your turn}
\begin{itemize}
\raggedright
  \item \emph{Say it:} \emph{kisha}
  \item \emph{Say it:} \emph{kisha}, once more
  \item \emph{Say it:} \emph{kisha ga keisatsu ni shitsumon shimashita} — A reporter put questions to the police.
  \item \emph{Recall:} say the Japanese for to print (in a paper), then the Japanese for an event, then say \emph{kisha} again
\end{itemize}

\begin{tcolorbox}[breakable,colback=teal!4,colframe=teal!35!black,title={Before you move on}]
What does \ja{きしゃ} mean? (\textquotedblleft{}A reporter; a journalist\textquotedblright{}.) Say it once more.
\end{tcolorbox}
\section[\texorpdfstring{\ja{くばる} (kubaru)}{くばる (kubaru)}]{\ja{くばる} (kubaru) — to hand out; to distribute}

\label{lesson:JA-C280-kubaru}



\begin{quote}
Before the new one: say the Japanese for an event, then the Japanese for a reporter.
\end{quote}

\subsection*{You'll want to know: \ja{くばる}}
\textbf{\ja{くばる}} — \emph{kubaru} — \textquotedblleft{}to hand out; to distribute\textquotedblright{}.

\begin{quote}
\textbf{\ja{えきで} \ja{しんぶんを} \ja{くばっていました。}} — \emph{eki de shinbun o kubatte imashita} — \textquotedblleft{}They were handing out newspapers at the station.\textquotedblright{}
\end{quote}

Polite \textbf{\ja{くばります}}, \emph{kubarimasu}; a godan verb. Past (polite): \textbf{\ja{こどもが} \ja{しんぶんを} \ja{くばりました。}}, \emph{kodomo ga shinbun o kubarimashita}.

\begin{cousinweb}
\textbf{\ja{くばる}} (\emph{kubaru}) is a native Japanese word (wago), inherited rather than borrowed from Chinese; its kanji was attached later as a kun'yomi (JMdict).
\end{cousinweb}

\begin{grammarlens}[title={What you've built}]
One more word for telling what happened.
\end{grammarlens}

\subsection*{Your turn}
\begin{itemize}
\raggedright
  \item \emph{Say it:} \emph{kubaru}
  \item \emph{Say it:} \emph{kubaru}, once more
  \item \emph{Say it:} \emph{eki de shinbun o kubatte imashita} — They were handing out newspapers at the station.
  \item \emph{Recall:} say the Japanese for an event, then the Japanese for a reporter, then say \emph{kubaru} again
\end{itemize}

\begin{tcolorbox}[breakable,colback=teal!4,colframe=teal!35!black,title={Before you move on}]
What does \ja{くばる} mean? (\textquotedblleft{}To hand out; to distribute\textquotedblright{}.) Say it once more.
\end{tcolorbox}
